\documentclass[10pt,aspectratio=169]{beamer}
\usetheme{Madrid}
\usecolortheme{seahorse}
\usepackage[utf8]{inputenc}
\usepackage[T1]{fontenc}
\usepackage{booktabs}
\usepackage{xcolor}

\title{Chern Calibration Copilot}
\subtitle{Diagnóstico de calibración cuántica: cómo el ruido degrada el número de Chern}
\author{Jesus · Hackathon AMD ACT III · Pista 1: Intelligent Industry}
\institute{AMD Developer Hackathon ACT III — lablab.ai × AMD}
\date{17 de octubre de 2026}

\begin{document}

\frame{\titlepage}

\section{Contexto}
\begin{frame}{El problema}
\begin{itemize}
\item El \emph{número de Chern} es un invariante topológico que clasifica fases de materia cuántica.
\item Cuando preparamos el estado cuántico con ruido, ¿cuánta infidelidad tolera antes de dejar de ser el valor exacto?
\item No hay respuesta publicada. Los ordenadores cuánticos actuales operan con fidelidades del 95--99\%, y no se sabe cuándo el invariante se rompe.
\end{itemize}
\end{frame}

\begin{frame}{Lo que hicimos}
\begin{itemize}
\item Simulador statevector en PyTorch/ROCm, validado contra PennyLane con error $10^{-15}$.
\item Preparamos el estado de la banda ocupada del modelo de Haldane con ruido depolarizante por puerta.
\item Medimos el Chern con el método Fukui--Hatsugai--Suzuki (FHS).
\item Probamos tres modelos de ruido con 256 semillas y barras de error.
\end{itemize}
\end{frame}

\section{Resultado}
\begin{frame}{Resultado principal}
\begin{table}[!htbp]
\centering
\small
\begin{tabular}{lcc}
\toprule
Modelo de ruido & $1-F$ & \% acierto (Chern exacto) \\
\midrule
\texttt{final} --- control coherente & 0,42 & \textbf{100 \%} \\
\texttt{shared} --- Pauli fijo por puerta & 0,42 & 28 \% \\
\texttt{indep} --- ruido por (puerta, k) & 0,42 & 0 \% \\
\bottomrule
\end{tabular}
\end{table}
\vspace{0.5cm}
\begin{block}{Hallazgo}
El invariante es \textbf{exactamente robusto} ante error coherente, se degrada suavemente ante error en posiciones fijas de puerta, y se destruye ante ruido incoherente. No hay un umbral único: depende de la estructura del error, no solo de su tamaño.
\end{block}
\end{frame}

\begin{frame}{El agente de calibración}
\begin{itemize}
\item Un agente LLM (Gemma vía AMD Token Factory) lee las especificaciones externas del dispositivo (presupuesto de error, coherencia, criterio de aceptación).
\item Emite un plan de mitigación \textbf{con fuentes citadas}.
\item El agente \emph{no} calcula el Chern: interpreta resultados.
\end{itemize}
\end{frame}

\section{AMD y honestidad}
\begin{frame}{AMD en el producto}
\begin{itemize}
\item Simulación en GPU AMD Radeon Pro W7900D (ROCm 7.15, 48 GiB).
\item Justificada por memoria: 256 semillas $\times$ n=14 $\times$ malla 12$\times$12 requieren $\sim$45 GiB, que superan los 6 GiB de CPU.
\item No se afirma speedup: no se ha medido en GPU localmente.
\end{itemize}
\end{frame}

\begin{frame}{Limitaciones declaradas}
\begin{itemize}
\item $n=1$ da el mismo Chern que $n=14$: el estado de Haldane es 1 cúbit por $k$.
\item No simulamos el efecto Hall cuántico físico: calculamos un invariante.
\item Los valores del corpus del agente son ilustrativos, no mediciones de un dispositivo real.
\end{itemize}
\end{frame}

\section{Demo}
\begin{frame}{Demo}
\begin{itemize}
\item \texttt{streamlit run demo/app.py} — barra de ruido, modelo de ruido, semillas, curva de acierto y plan del agente.
\item Repo: \url{github.com/jalp17/chern-calibration-copilot}
\end{itemize}
\end{frame}

\begin{frame}{Referencias}
\begin{itemize}
\item Lado et al., \emph{Phys. Rev. Research} 6, 043288 (2024), arXiv:2404.06048.
\item Fukui, Hatsugai, Suzuki, \emph{JPSJ} 74, 1674 (2005).
\item Haldane, \emph{PRL} 61, 2015 (1988).
\end{itemize}
\end{frame}

\end{document}
